\chapter{Methodology}
\label{chap:methodology}

This chapter describes how the study was carried out: the dataset, the
preprocessing and feature engineering, the four model architectures, the
training protocol, and the metrics used for evaluation, data efficiency and
calibration. The complete implementation is released as open
source.\footnote{\url{https://github.com/ramc77/fyp-jet-tagging}}

\section{Dataset}
\label{sec:dataset}

We use the public \emph{Top-Quark Tagging Reference
Dataset}~\citep{kasieczka2019dataset}, the community-standard benchmark for
this problem~\citep{kasieczka2019landscape}. It contains $2\times10^{6}$
simulated jets from $14$\,TeV proton--proton collisions, generated with
\textsc{Pythia\,8}~\citep{sjostrand2015pythia} and a \textsc{Delphes}
fast detector simulation~\citep{delphes2014} of an ATLAS-like detector.
Signal jets come from boosted hadronic top decays; background jets come from
QCD dijet production. Jets are clustered with the anti-$k_t$
algorithm~\citep{cacciari2008} with radius $R=0.8$ and transverse momentum
$550<p_T<650$\,GeV, and each jet is stored as up to $200$ constituent
four-momenta with a binary label (top vs.\ QCD).

For computational tractability we draw stratified subsamples of up to
$5\times10^{4}$ jets per split (training, validation, test), keeping the
roughly $50/50$ signal-to-background balance of the full dataset. The deep
models use the leading $100$ constituents per jet, ordered by $p_T$.

\section{Preprocessing and Feature Engineering}
\label{sec:features}

\subsection{Per-constituent features}
From each constituent four-momentum $(E,p_x,p_y,p_z)$ we compute the
transverse momentum $p_T$, pseudorapidity $\eta$ and azimuth $\phi$, and the
coordinates relative to the $p_T$-weighted jet axis,
$(\Delta\eta,\Delta\phi)$. The resulting seven-dimensional per-constituent
feature vector
\begin{equation}
\mathbf{x}=\bigl(p_T,\;\eta,\;\phi,\;E,\;\Delta\eta,\;\Delta\phi,\;\log p_T\bigr)
\end{equation}
is standardised feature-wise using statistics computed on the training split
only.

\subsection{Engineered jet-level features (for the BDT)}
For the BDT baseline we compute thirteen high-level observables: the jet
mass, $p_T$ and $\eta$; the constituent multiplicity; the jet width; the
leading and sub-leading $p_T$ fractions; the $N$-subjettiness variables
$\tau_1,\tau_2,\tau_3$ and ratios $\tau_{21},\tau_{32}$~\citep{thaler2011};
and the energy-correlation ratio $C_2$~\citep{larkoski2013}. Their
distributions for signal and background are shown in
Chapter~\ref{chap:results} and behave exactly as physics expects, validating
the pipeline.

\subsection{Jet images (for the CNN)}
For the CNN, constituents are binned into $40\times40$ images in the
$(\Delta\eta,\Delta\phi)$ plane with three channels: the $p_T$ sum, the
constituent multiplicity, and the $p_T^2$ sum per pixel. Each channel is
normalised by its per-image maximum.

\subsection{Pairwise interaction features (for the Transformer)}
For the Particle Transformer we additionally compute, on the fly, four
pairwise quantities for every constituent pair $(i,j)$:
$\ln\Delta R_{ij}$, $\ln k_{T,ij}$, the momentum fraction $z_{ij}$, and
$\Delta\eta_{ij}$, following~\citet{qu2022particletransformer}.

\section{Model Architectures}
\label{sec:models}

\subsection{Boosted decision tree (BDT)}
A gradient-boosted decision tree trained with
\textsc{XGBoost}~\citep{chen2016xgboost} on the thirteen engineered features,
with up to $500$ boosting rounds, maximum depth $7$, learning rate $0.1$ and
column/row subsampling $0.8$, with early stopping on validation AUC. It is
the interpretable, classical baseline.

\subsection{Convolutional neural network (CNN)}
A compact ResNet-style network~\citep{he2016resnet} with three residual
blocks ($64\rightarrow128\rightarrow256$ channels), batch normalisation,
global average pooling and a fully connected head, acting on the
$40\times40\times3$ jet images.

\subsection{ParticleNet}
A from-scratch implementation of
ParticleNet~\citep{qu2020particlenet}, using three EdgeConv
blocks~\citep{wang2019dgcnn} with a $k=16$ dynamic nearest-neighbour graph
in $(\Delta\eta,\Delta\phi)$ space, followed by global pooling and a fully
connected classifier. It treats the jet as a permutation-invariant particle
cloud.

\subsection{Particle Transformer}
A self-attention model~\citep{qu2022particletransformer} with three encoder
blocks, multi-head attention, and the four pairwise features injected as an
additive attention bias. A learnable class token aggregates the sequence for
classification, and its attention weights provide an interpretability handle.

\section{Training Protocol}
\label{sec:training}

All deep models are trained by minimising binary cross-entropy with the
Adam/AdamW optimiser~\citep{kingma2014adam,loshchilov2019adamw}, with learning
rates of $10^{-3}$--$10^{-4}$, a plateau learning-rate scheduler, gradient
clipping and early stopping on validation AUC. Training was performed on a
free Google Colaboratory T4 GPU; the same code runs on CPU. To make long runs
robust, the pipeline checkpoints after every epoch and resumes automatically,
so a session can be interrupted and continued without loss of progress.

\section{Evaluation Metrics}
\label{sec:metrics}

\subsection{Discrimination}
We report the area under the ROC curve (AUC), the accuracy at a score
threshold of $0.5$, and the \emph{background rejection}
$1/\varepsilon_B$ at fixed signal efficiencies $\varepsilon_S=0.5$ and
$0.3$, which is the figure of merit most relevant to physics analyses.
Uncertainties on the AUC are estimated by bootstrap resampling of the test
set.

\subsection{Data efficiency}
Each model is retrained from scratch at
$N_{\text{train}}\in\{5{,}000,\,10{,}000,\,25{,}000,\,50{,}000\}$ jets, with
all other settings fixed, and the resulting test AUC and background rejection
are plotted against $N_{\text{train}}$ to expose how each architecture scales.

\subsection{Calibration}
We measure the Expected Calibration Error
\begin{equation}
\mathrm{ECE}=\sum_{m=1}^{M}\frac{|B_m|}{N}\,
\bigl|\,\mathrm{acc}(B_m)-\mathrm{conf}(B_m)\,\bigr|,
\label{eq:ece}
\end{equation}
where the predictions are partitioned into $M$ equal-width score bins $B_m$
and $\mathrm{acc}$ and $\mathrm{conf}$ are the bin accuracy and mean
confidence~\citep{guo2017calibration}. We complement the ECE with reliability
diagrams, and we fit a single \emph{temperature} $T$ on the validation set by
minimising the negative log-likelihood of $\sigma(\,z/T\,)$, where $z$ are
the logits, reporting the ECE before and after this rescaling.
